La presència de l'isòtop ferro 60 ($^{60}$Fe) en algunes roques lunars i en alguns sediments oceànics indica, segons alguns astrofísics, que una supernova va esclatar a les proximitats del Sistema Solar en una època relativament recent (a escala còsmica) i va fer arribar aquest isòtop fins a la Terra. El $^{60}$Fe té un període de semidesintegració de 2,6 milions d'anys.

\begin{apartat}{1}
Si hi hagués hagut $^{60}$Fe quan la Terra es va formar, fa 4\,400 milions d'anys, quin percentatge d'aquest $^{60}$Fe primordial quedaria ara? Si el $^{60}$Fe es va originar en l'explosió d'una supernova fa 13 milions d'anys, quin percentatge d'aquest $^{60}$Fe hauria de quedar encara?
\end{apartat}

\begin{apartat}{1}
El $^{60}$Fe es transforma, mitjançant una desintegració $\beta^-$, en un isòtop de cobalt (Co) de vida breu, el qual torna a patir una nova desintegració $\beta^-$ i produeix un isòtop estable de níquel (Ni). Escriviu les equacions nuclears de les dues desintegracions, incloent-hi els antineutrins.
\end{apartat}

\blocetiquetat{Dada:}{Nombre atòmic del ferro (Fe): 26.}
